\begin{table}[htbp]
  \centering
  \small
  \caption{Description textuelle du cas d'utilisation : «~Consulter le tableau de bord corporate~»}
  \label{tab:cu-tableau-de-bord}
  \PFETableSetup
  \PFETableRowColors
  \PFETableArrayStretch
  \PFETableTabularXBegin
  \begin{tabularx}{\PFETableBlockWidth}{@{}>{\raggedright\arraybackslash}p{3.2cm} >{\raggedright\arraybackslash}X@{}}
    \tableHeaderStyle \tableHeaderCell{Champ} & \tableHeaderCell{Contenu} \\
    \textbf{Titre} & Consulter le tableau de bord corporate \\
    \textbf{Acteurs} & Utilisateur corporate (pilotage), éventuellement rôles restreints en lecture. \\
    \textbf{Objectif} &
      Offrir une vue consolidée des plans, activités et indicateurs clés pour suivre l'avancement CSR multi-sites et détecter les écarts. \\
    \textbf{Préconditions} &
      L'utilisateur est authentifié avec un périmètre corporate. Les données sources (plans validés, activités, réalisations) sont disponibles. \\
    \textbf{Postconditions} &
      Les indicateurs affichés reflètent les filtres choisis ; l'utilisateur peut basculer vers le détail d'un site ou exporter. \\
    \textbf{Scénario nominal} &
      L'utilisateur ouvre le tableau de bord, applique des filtres (année, site, catégorie), consulte les graphiques et tableaux, puis approfondit une zone sensible (top activités, chronologie, performance site). \\
    \textbf{Scénario alternatif} &
      \textbf{Données partielles :} affichage d'un bandeau d'avertissement si certains sites n'ont pas encore soumis de plan.\par\medskip
      \textbf{Timeout :} message invitant à réduire le périmètre temporel ou à réessayer. \\
  \end{tabularx}%
  \PFETableTabularXEnd
\end{table}

\begin{table}[htbp]
  \centering
  \small
  \caption{Description textuelle du cas d'utilisation : «~Exploiter les données pour le reporting et Power BI~»}
  \label{tab:cu-reporting-powerbi}
  \PFETableSetup
  \PFETableRowColors
  \PFETableArrayStretch
  \PFETableTabularXBegin
  \begin{tabularx}{\PFETableBlockWidth}{@{}>{\raggedright\arraybackslash}p{3.2cm} >{\raggedright\arraybackslash}X@{}}
    \tableHeaderStyle \tableHeaderCell{Champ} & \tableHeaderCell{Contenu} \\
    \textbf{Titre} & Exploiter les données pour le reporting et Power BI \\
    \textbf{Acteurs} & Système (synchronisation), utilisateur corporate (consommation des rapports). \\
    \textbf{Objectif} &
      Permettre l'alimentation de tableaux de bord Power BI et la réutilisation des agrégats CSR dans des analyses externes ou des revues périodiques. \\
    \textbf{Préconditions} &
      Jeux de données identifiés ; comptes ou clés d'accès Power BI conformes à la politique de sécurité de l'entreprise. \\
    \textbf{Postconditions} &
      Les rapports distants reflètent une photographie cohérente des données ; la fréquence de rafraîchissement est documentée. \\
    \textbf{Scénario nominal} &
      Le système exporte ou publie les agrégats vers le canal prévu ; l'utilisateur corporate ouvre le rapport Power BI et croise les vues avec d'autres sources internes si besoin. \\
    \textbf{Scénario alternatif} &
      \textbf{API Power BI en cours de complétion :} usage intermédiaire d'exports fichiers depuis le tableau de bord web.\par\medskip
      \textbf{Échec de synchronisation :} journal d'erreur et notification aux administrateurs. \\
  \end{tabularx}%
  \PFETableTabularXEnd
\end{table}
